% 第九卷 13.1.3 节：RAG与FiD接入多个检索片段的位置（跨栏，约150 mm宽）。
% 蓝色：输入与片段；黄色描边：模型部件；红色：输出概率或答案。
\begin{tikzpicture}[
  x=1mm,y=1mm,
  figure picture,
  inp/.style={figure node,draw=FigureInput,fill=FigureInputFill,
    minimum height=7mm,minimum width=18mm,inner xsep=2pt,inner ysep=2pt,
    font=\figurefont\scriptsize},
  mod/.style={figure node,draw=FigureModel,fill=FigureModelFill,
    minimum height=7mm,minimum width=18mm,inner xsep=2pt,inner ysep=2pt,
    font=\figurefont\scriptsize},
  res/.style={figure node,draw=FigureOutput,fill=FigureOutputFill,
    minimum height=8mm,minimum width=60mm,inner xsep=2pt,inner ysep=2pt,
    font=\figurefont\scriptsize},
  note/.style={font=\figurefont\scriptsize,text=FigureMuted,align=center},
  fl/.style={figure flow}
]
  \figurePanel{0}{0}{72}{86}{RAG：在概率层对片段边缘化}{FigureInputFill}
  \figurePanel{78}{0}{72}{86}{FiD：在解码器内融合片段}{FigureInputFill}

  % ---------- RAG ----------
  \node[inp,minimum width=24mm] (rq) at (36,8) {问题$x$};
  \node[mod,minimum width=60mm] (rr) at (36,20)
    {DPR检索器：$p_\eta(z\mid x)\propto\exp(\bm d(z)^{\top}\bm q(x))$，取前$k$个};
  \foreach \i/\xx/\zz in {1/14/{z_1},2/36/{z_2},3/58/{z_k}}{
    \node[inp] (rz\i) at (\xx,32) {片段$\zz$};
    \node[mod] (rg\i) at (\xx,46) {生成器$p_\theta$\\条件$x,\zz$};
    \draw[fl] (rz\i) -- (rg\i);
    \draw[fl] (rr.north -| rz\i) -- (rz\i);
  }
  \node[note] at (47,32) {$\cdots$};
  \node[res] (ro) at (36,62)
    {$\sum_{z}p_\eta(z\mid x)\times$（整句概率或每个词元概率）};
  \foreach \i in {1,2,3}{\draw[fl] (rg\i) -- (rg\i |- ro.south);}
  \draw[fl] (rq) -- (rr);
  \node[note] at (36,73) {每个分支只读一个片段；\\检索概率作为权重参与求和};

  % ---------- FiD ----------
  \foreach \i/\xx/\zz in {1/92/{z_1},2/114/{z_2},3/136/{z_k}}{
    \node[inp] (fz\i) at (\xx,8) {$x$与片段$\zz$};
    \node[mod] (fe\i) at (\xx,21) {编码器\\$\bm H_{\i}\in\mathbb R^{\ell\times d}$};
    \draw[fl] (fz\i) -- (fe\i);
  }
  \node[note] at (125,8) {$\cdots$};
  \node[mod,minimum width=60mm] (fc) at (114,35)
    {拼接$[\bm H_1;\ldots;\bm H_k]\in\mathbb R^{k\ell\times d}$};
  \foreach \i in {1,2,3}{\draw[fl] (fe\i) -- (fe\i |- fc.south);}
  \node[mod,minimum width=60mm] (fd) at (114,49)
    {单个解码器：每步交叉注意全部$k\ell$个位置};
  \node[res] (fo) at (114,62) {$\prod_i p_\theta(y_i\mid y_{<i},[\bm H_1;\ldots;\bm H_k])$};
  \draw[fl] (fc) -- (fd);
  \draw[fl] (fd) -- (fo);
  \node[note] at (114,73) {编码代价随$k$线性增长；\\一个词元可同时依据多个片段};
\end{tikzpicture}
